\section[Maths]{Maths}

\blockframe{3}{Maths}{The reason most people are here}

\begin{frame}
    \frametitle{Where we are going}
    \vfill
    \begin{roadmap}{3}
        \roadmapresults
    \end{roadmap}
    \vfill
\end{frame}

\begin{frame}
    \frametitle{Warm-up: what comes out?}
    \begin{center}
        \tcbinputlst[][width=.6\textwidth]{listings/03-warmup.tex}
    \end{center}
    \vfill
    \begin{alertblock}<2->{}
        No empty line: \textbf{not} a new paragraph
    \end{alertblock}
    \note{%
        \textbf{Two minutes, in pairs, laptops shut.}
        \begin{itemize}
            \item Nasty on purpose.
                Those who get it right will be pleased with themselves for the rest of the hour;
                those who get it wrong made the mistake where it costs nothing.
            \item Tie it back: ``the empty line is the paragraph break.''
                Same rule as before the break, in a place where it is easy to forget.
        \end{itemize}
    }
\end{frame}

\begin{frame}
    \frametitle{Two ways in}
    \begin{columns}[T, onlytextwidth]
        \begin{column}{.48\textwidth}
            {\large\bfseries Inline, inside a sentence\par}
            \medskip
            \tcbinputlst[fontsize=\footnotesize]{listings/03-inline.tex}
            \begin{outputbox}
                The constant $\tau$ sets the pace,
and $T_0$ is where we start.

            \end{outputbox}
        \end{column}
        \begin{column}{.48\textwidth}
            {\large\bfseries Display, on its own line\par}
            \medskip
            \tcbinputlst[fontsize=\footnotesize]{listings/03-display.tex}
            \begin{outputbox}
                \[
    T(t) = T_\mathrm{env} + \Delta T
\]

            \end{outputbox}
        \end{column}
    \end{columns}
    \vfill
    \begin{alertblock}<2->{}
        \nope{} \texttt{\$\$...\$\$}
        \qquad
        \yep{} \texttt{\textbackslash[...\textbackslash]}
    \end{alertblock}
    \note{%
        \begin{itemize}
            \item \texttt{\$\$...\$\$} is a plain-\TeX{} leftover and breaks the spacing.
            \item The \texttt{\$\$} warning matters: they \emph{will} have met it in Markdown, Jupyter
                and every StackExchange answer.
                Say it once, clearly; the handout repeats it.
            \item Inline maths shrinks to fit the line (small fractions, compact limits), display maths does not.
                Not a bug: it keeps the line spacing even.
        \end{itemize}
    }
\end{frame}

% Three states: unnumbered, numbered, labelled and referenced.
% Three separate frames rather than overlays: beamer only sets a \label on
% the first overlay of a frame, so the label would never exist otherwise.
\begin{frame}
    \frametitle{Numbered, and worth pointing at}
    \tcbinputlst{listings/03-equation-1.tex}
    \begin{outputbox}
        \[
    T(t) = T_\mathrm{env} + (T_0 - T_\mathrm{env}) \mathrm{e}^{-t/\tau}
\]

    \end{outputbox}
    \note{%
        Three steps: unnumbered, numbered, labelled and referenced.
    }
\end{frame}

\begin{frame}
    \frametitle{Numbered, and worth pointing at}
    \tcbinputlst[highlightlines={1,3}]{listings/03-equation-2.tex}
    \begin{outputbox}
        \begin{equation}
    T(t) = T_{\mathrm{env}} + (T_0 - T_{\mathrm{env}}) \mathrm{e}^{-t/\tau}
\end{equation}

    \end{outputbox}
\end{frame}

\begin{frame}
    \frametitle{Numbered, and worth pointing at}
    \tcbinputlst[highlightlines={3,5}]{listings/03-equation-3.tex}
    \begin{outputbox}
        \begin{equation}
    T(t) = T_\mathrm{env} + (T_0 - T_\mathrm{env}) \mathrm{e}^{-t/\tau}
    \label{eq:cooling}
\end{equation}
As Equation~\eqref{eq:cooling} shows, the drink cools fastest at the start.

    \end{outputbox}
    \vfill
    \begin{exampleblock}<2->{}
        The number is never typed
    \end{exampleblock}
    \note{%
        \begin{itemize}
            \item Add an equation above it and everything renumbers.
            \item The first callback to the block 1 promise.
                Name it: ``remember the contents list? Same deal.''
            \item \texttt{\textbackslash eqref} gives ``(1)'' with the brackets, \texttt{\textbackslash ref} a bare ``1''.
                For equations you almost always want \texttt{\textbackslash eqref}.
        \end{itemize}
    }
\end{frame}

\begin{frame}
    \frametitle{Several lines, one alignment}
    \only<1>{\tcbinputlst[fontsize=\footnotesize]{listings/03-align.tex}}%
    \only<2>{\tcbinputlst[fontsize=\footnotesize, highlightlines={3,4}]{listings/03-align.tex}}%
    \only<3->{\tcbinputlst[fontsize=\footnotesize, highlightlines={3}]{listings/03-align.tex}}%
    \begin{outputbox}
        \begin{align}
    T(t) &= T_{\mathrm{env}} + \Delta T \cdot \mathrm{e}^{-t/\tau} \\
    \Delta T &= T_0 - T_{\mathrm{env}} \notag
\end{align}

    \end{outputbox}
    \begin{columns}[T, onlytextwidth]
        \begin{column}{.47\textwidth}
            \uncover<2->{%
                \texttt{\&}: align here, before each \texttt{=}
            }
        \end{column}
        \begin{column}{.47\textwidth}
            \uncover<3->{%
                \texttt{\textbackslash\textbackslash}: new line\par
                \texttt{\textbackslash notag}: no number for this line
            }
        \end{column}
    \end{columns}
    \note{%
        \begin{itemize}
            \item \texttt{\textbackslash notag} is for a line nobody will cite.
            \item Callback: in block 2 you said \texttt{\textbackslash\textbackslash} belongs in two places,
                tables and maths.
                This is one of them.
                Say ``told you''; the payoff for the earlier promise makes the next one credible.
            \item Do not teach \texttt{align*}, \texttt{gather}, \texttt{split}, \texttt{multline} or \texttt{aligned}.
                \texttt{align} plus \texttt{\textbackslash notag} covers everything they need this year;
                the handout lists the rest.
        \end{itemize}
    }
\end{frame}

\begin{frame}
    \frametitle{Labels are not variables}
    \begin{columns}[T, onlytextwidth]
        \begin{column}{.48\textwidth}
            {\large\bfseries\nope{} Wrong\par}
            \smallskip
            \tcbinputlst[fontsize=\footnotesize]{listings/03-subscript-wrong.tex}
            \begin{outputbox}
                $V_{out}$

                \hfill
                {\small\color{black!60} $V \cdot o \cdot u \cdot t$?}
            \end{outputbox}
        \end{column}
        \begin{column}{.48\textwidth}
            {\large\bfseries\yep{} Right\par}
            \smallskip
            \tcbinputlst[fontsize=\footnotesize]{listings/03-subscript-right.tex}
            \begin{outputbox}
                $V_\mathrm{out}$

                \hfill
                {\small\color{black!60} upright}
            \end{outputbox}
        \end{column}
    \end{columns}
    \vfill
    \begin{exampleblock}<2->{}
        Variables italic, labels upright: \texttt{\textbackslash mathrm}
    \end{exampleblock}
    \note{%
        \begin{itemize}
            \item \texttt{V\_\{out\}} is the most common thing beginners get wrong in a report,
                and it is visible from the back of the room.
            \item ``out'' is a label, not three variables multiplied together.
        \end{itemize}
    }
\end{frame}

\begin{frame}
    \frametitle{Words inside formulas}
    \tcbinputlst{listings/03-cases.tex}
    % cases only works in maths mode, hence the \[ \] around it
    \begin{outputbox}
        \[
            \[
    \begin{cases}
        \text{too hot} & \text{if } T > 65\,^{\circ}\mathrm{C} \\
        \text{drinkable} & \text{otherwise}
    \end{cases}
\]

        \]
    \end{outputbox}
    \vfill
    \begin{exampleblock}<2->{}
        Words in a formula: \texttt{\textbackslash text\{...\}}
    \end{exampleblock}
    \note{%
        \begin{itemize}
            \item A quantity goes in maths mode, a word goes in \texttt{\textbackslash text}, even mid-sentence.
            \item \texttt{cases} shows \texttt{\&} and \texttt{\textbackslash\textbackslash} doing the same job as in \texttt{align}.
                Point that out.
            \item The template's version has a third branch: same shape, one more line.
            \item They will find \texttt{\textbackslash,\^{}\{\textbackslash circ\}\textbackslash mathrm\{C\}} ugly.
                Good: tell them a better way is coming at the end.
                That is the hook for \texttt{siunitx}.
        \end{itemize}
    }
\end{frame}

\begin{frame}[fragile]
    \frametitle{Dos and don'ts}
    \begin{center}
        \renewcommand{\arraystretch}{1.25}
        \begin{tabular}{@{}lll@{}}
            \toprule
            \textbf{Don't}     & \textbf{Do}                                  & \textbf{Why} \\
            \midrule
            \code|a * b|       & \code|a \cdot b|                             & \texttt{*} is convolution \\
            \code|$$...$$|     & \code|\[...\]|                               & plain-\TeX{} leftover \\
            \code|V_{out}|     & \code|V_\mathrm{out}|                        & a label, not variables \\
            \code|sin(x)|      & \code|\sin(x)|                               & upright, spaced \\
            \code|d/dx|        & \code|\frac{\mathrm{d}f}{\mathrm{d}x}|       & operator, upright \\
            \bottomrule
        \end{tabular}
    \end{center}
    \vfill
    \begin{exampleblock}<2->{}
        \code|\cdot|: product
        \qquad
        \code|\times|: cross product, dimensions
    \end{exampleblock}
    \note{%
        \begin{itemize}
            \item Two minutes.
                It is a reference table and it is in the handout;
                reading it out line by line will lose the room.
            \item Pick the two that matter for \emph{this} group.
                For a mixed audience: \texttt{\textbackslash mathrm} for labels and \texttt{\textbackslash sin} for functions.
            \item The 2025 deck said ``never use \texttt{\textbackslash times}''.
                Too strong: it is right for cross products and dimensions,
                and someone in the room will know that.
        \end{itemize}
    }
\end{frame}

\begin{frame}
    \frametitle{Don't know the symbol?}
    \vfill
    \begin{center}
        {\LARGE\bfseries\href{https://detexify.kirelabs.org}{detexify.kirelabs.org}\par}
        \bigskip
        {\Large Draw a symbol, get the command\par}
    \end{center}
    \vfill
    \note{%
        Ten seconds, and worth every one.
        Nobody memorises the symbol tables.
        \begin{itemize}
            \item The fear of memorising hundreds of commands is the main reason people bounce off \LaTeX{}.
                This kills it.
            \item The same is true of everything today: you look it up until you stop needing to.
        \end{itemize}
    }
\end{frame}

\begin{frame}
    \frametitle{Your turn}
    \handson{3}{3.1--3.6}{5--10}{checkpoint-03.tex}{%
        Fill in the Theory section.\par
        \medskip
        {\normalsize\color{black!60}
        load \texttt{amsmath} $\rightarrow$ set the symbol list in maths
        $\rightarrow$ a numbered equation with a label $\rightarrow$ an \texttt{align} derivation
        $\rightarrow$ refer back to both $\rightarrow$ a \texttt{cases} block}
    }
    \note{%
        \textbf{5 to 10 minutes.}
        The fullest block, so expect fewer people to finish.
        That is planned for; it is what checkpoint-03 is for.
        \begin{itemize}
            \item Anyone who finishes: the three drills at the bottom of their file (they come with checkpoint-02)
                (logistic map, Shannon entropy, inverse Laplace transform).
                The Laplace one keeps the strongest person busy until you stop.
            \item Walk the room looking for \emph{one} thing:
                missing \texttt{\$} around symbols in prose.
            \item \textbf{Break after this}, 5 minutes.
                Last one.
        \end{itemize}
    }
\end{frame}
